\documentclass[journal]{IEEEtran}

\usepackage[utf8]{inputenc}
\usepackage[T1]{fontenc}
\usepackage{amsmath,amssymb,amsfonts}
\usepackage{graphicx}
\usepackage{booktabs}
\usepackage{array}
\usepackage{url}
\usepackage[hidelinks]{hyperref}
\usepackage{orcidlink}

\begin{document}

\title{Physics-Informed Reconstruction of Unmeasurable Cardiac Fields from Real
Data: Electrocardiographic Imaging and 4D-Flow Aortic Pressure, with Honest Nulls}

\author{%
  Felipe Santib\'a\~nez-Leal~\orcidlink{0000-0002-0150-3246},~\IEEEmembership{Independent Researcher}%
  \thanks{F.\ Santib\'a\~nez-Leal is with the CAOS open-research programme,
    Santiago, Chile (e-mail: \texttt{fsantibanezleal@ug.uchile.cl};
    ORCID: 0000-0002-0150-3246).}%
  \thanks{Preprint, 2026. Code, derived reconstruction artifacts, and the
    interactive viewer are released under MIT at
    \url{https://github.com/fsantibanezleal/CAOS_RES_CardioPINN}; live at
    \url{https://cardiopinn.fasl-work.com}. This is a validated methodological
    result on real experimental data, not a clinically deployed system; raw
    datasets are used under their data-use agreements and are not redistributed.}%
}

\markboth{Preprint, 2026 --- CAOS open-research programme}{Santib\'a\~nez-Leal:
Physics-Informed Reconstruction of Unmeasurable Cardiac Fields from Real Data}

\IEEEtitleabstractindextext{%
\begin{abstract}
Two clinically decisive cardiac quantities cannot be measured where they are
needed: the heart-surface potential map (measured only with an invasive cage) and
the intra-vascular pressure field (measured only with a catheter). Both are tied by
a governing equation to a quantity that \emph{can} be measured non-invasively,
the body-surface potential and the 4D-flow velocity, so both can be reconstructed.
We report a real-data-only study of this reconstruction across two distinct physics
domains, with a strict discipline: every case fits a real measured signal and is
validated against a real gold standard or an exact analytic gate; there is no
synthetic ground truth, because a network that only re-solves an equation a
classical solver already solves answers no clinical question. \emph{Case 1
(electrocardiographic imaging, quasi-static volume conduction).} On real EDGAR
experiments reconstructed by an identical pipeline with no per-heart retuning, a
forward operator on the real geometry with a graph-Laplacian prior and a
deep-ensemble node uncertainty recovers heart-surface potentials at a relative
error of 0.54--0.65 and spatial correlation 0.72--0.85 against the measured cage,
with $\approx 90\%$ of nodes inside two standard deviations, and paced beats
reconstructing better than sinus as physically expected. A full boundary-element
forward operator, analytic-gated to correlation 1.00 on a concentric-sphere
problem, honestly does \emph{not} beat the calibrated single-layer on the real
coarse electrode geometry (dog: 0.54 vs 0.63 relative error); the null is reported.
\emph{Case 2 (4D-flow aortic pressure, incompressible Navier--Stokes).} A
divergence-free velocity network denoises a real thoracic-aorta 4D-flow scan (raw
divergence reduced $2.3\times$) and the relative pressure is forced out by the
pressure-Poisson equation with the source and Neumann flux computed from the
network's analytic derivatives rather than by finite differences at the lumen edge,
which is what takes the recovered pressure from a non-physiological thousands of
mmHg to a physiological 0.79 mmHg range for this unobstructed aorta, the same order
as the clinical simplified-Bernoulli estimate (2.51 mmHg). The solve is analytic-gated
to correlation 1.00 on a converging duct; the momentum-residual PINN that leaves
pressure gauge-free is kept as the documented failed baseline. The common thread is
methodological: separate the well-posed part (velocity, potentials) from the
ill-posed part (pressure, inverse source), validate against a real gold standard or
an exact gate, and publish the nulls.
\end{abstract}

\begin{IEEEkeywords}
physics-informed neural networks, electrocardiographic imaging, inverse problems,
4D-flow MRI, pressure-Poisson equation, Navier--Stokes, volume conduction,
boundary element method, deep ensembles, uncertainty quantification, cardiac
reconstruction.
\end{IEEEkeywords}}

\maketitle
\IEEEdisplaynotcompsoctitleabstractindextext
\IEEEpeerreviewmaketitle

% =====================================================================
\section{Introduction}
% =====================================================================

\IEEEPARstart{P}{hysics-informed} neural networks (PINNs) embed a governing
differential equation in the training loss, so a network can fit sparse or noisy
measurements while respecting the physics that relates them to unobserved
fields~\cite{raissi2019}. Their most compelling use is not re-solving a forward
problem a classical solver already handles, but recovering a field that
\emph{cannot} be measured from one that can~\cite{raissi2020}. Two cardiac problems
have exactly this shape. In electrocardiographic imaging (ECGi), the clinician needs
the heart-surface potential map to localize an arrhythmia, but a patient offers only
the body-surface recording; the true heart-surface map, from an electrode cage, is a
gold standard one never has in the clinic~\cite{ramanathan2004}. In vascular flow,
the clinician needs the pressure drop across a stenosis, but 4D-flow MRI measures
only velocity; pressure is obtained invasively by catheter~\cite{krittian2012}. In
both, a governing equation ties the wanted field to the measured one.

This paper is deliberately real-data-only. A recurring failure mode in
physics-informed reconstruction is to validate on synthetic data, where a network is
scored against a field a solver produced under the same assumptions the network
uses; such a result demonstrates code correctness, not a clinical capability. We
therefore fit only real measured signals and validate only against a real measured
gold standard or, where no gold standard can exist (pressure is never measured
non-invasively), against an exact analytic gate that must pass before any real datum
is trusted. We report two cases in two different physics domains and, crucially, the
honest negative results each produced: a boundary-element forward operator that does
not beat a simpler one on real geometry, and a momentum-residual PINN that cannot
recover pressure at aortic Reynolds number. Both nulls are informative and are kept
on the record.

% =====================================================================
\section{Case 1: Electrocardiographic imaging (volume conduction)}
\label{sec:ecgi}
% =====================================================================

\subsection{Problem and forward operator}

In the quasi-static approximation the torso is a volume conductor and the potential
$\phi$ obeys Laplace's equation $\nabla\!\cdot(\sigma\nabla\phi)=0$ between the
heart surface and the body surface. Discretized on the real geometry, the
body-surface potentials $b$ are a linear map of the heart-surface potentials $x$,
\begin{equation}
b = A\,x + \varepsilon,
\end{equation}
where $A$ is the transfer (lead-field) matrix and $\varepsilon$ is measurement
noise. Recovering $x$ from $b$ is severely ill-posed: the singular values of $A$
decay rapidly, so noise in $b$ is amplified without bound. We build $A$ two ways: a
single-layer point-source Green's-function approximation, and a full
boundary-element operator (BEM) whose double-layer entries use the exact
solid angle of a plane triangle~\cite{vanoosterom1983}.

\subsection{Regularized inversion with a spatial prior and node uncertainty}

The reconstruction is a Tikhonov problem with a graph-Laplacian prior $L$ on the
heart surface,
\begin{equation}
\hat{x} = \arg\min_x \; \lVert A\,x - b \rVert_2^2 + \lambda\,\lVert L\,x \rVert_2^2,
\label{eq:tikhonov}
\end{equation}
with $\lambda$ chosen on the L-curve~\cite{hansen1992}. Because a single regularized
point estimate carries no honest error bar, we draw an ensemble over realistic
measurement-noise perturbations of $b$ and re-solve, yielding a per-node predictive
distribution~\cite{lakshminarayanan2017}; the ensemble spread is recalibrated so that
its two-standard-deviation band covers the true per-node error at the nominal rate.

\subsection{Real validation on EDGAR}

The pipeline is applied, with no per-heart retuning, to independent real EDGAR
experiments~\cite{aras2015}: a human torso tank (192 body electrodes $\to$ 256-node
cage; sinus and two paced beats) and an in-situ dog (140 body $\to$ 1321-node
epicardium; sinus). The recovered heart-surface potentials are compared to the real
measured cage potentials with the standard relative error (RE) and spatial
correlation (CC)~\cite{cluitmans2018} (Table~\ref{tab:ecgi}). The numbers are
literature-consistent for torso-tank ECGi, the node uncertainty covers
$\approx 90\%$ of nodes at two standard deviations, and the paced beats reconstruct
better than sinus (higher CC), which is physically expected: a focal paced
activation is easier to localize than a diffuse sinus wavefront.

\begin{table}[!t]
\caption{Real ECGi validation against the measured cage potentials (identical
pipeline, no per-heart retuning). RE: relative error; CC: spatial correlation;
UQ: fraction of nodes within two standard deviations.}
\label{tab:ecgi}
\centering
\begin{tabular}{llrrr}
\toprule
Dataset & Beat & RE & CC & UQ ($2\sigma$) \\
\midrule
Human torso tank & Sinus & 0.65 & 0.72 & 0.90 \\
Human torso tank & Paced (PVP) & 0.58 & 0.80 & 0.89 \\
Human torso tank & AV-paced (AVP) & 0.54 & 0.85 & 0.90 \\
In-situ dog & Sinus & 0.54 & 0.78 & 0.90 \\
\bottomrule
\end{tabular}
\end{table}

\subsection{Honest null: BEM does not beat the single-layer here}

A more physically complete forward operator is not automatically a better one. The
BEM is analytic-gated on a concentric-sphere problem (correlation 1.00, error halving
per mesh refinement), so its implementation is correct. Yet on the real electrode
geometry it does not beat the calibrated single-layer: it requires closed 2-manifold
surfaces (the human torso-tank surface is open, so the BEM applies only to the dog
case), and where it applies the coarse 140-node torso makes the reconstruction
regularization-dominated (dog: single-layer RE 0.54 vs.\ BEM RE 0.63). Forward-operator
fidelity is not the bottleneck at this electrode density; the single-layer stays the
default, and the BEM is expected to matter only as electrode density and mesh closure
improve. This null is reported, not hidden.

% =====================================================================
\section{Case 2: 4D-flow aortic pressure (Navier--Stokes)}
\label{sec:flow}
% =====================================================================

\subsection{Problem and the pressure-Poisson route}

For an incompressible Newtonian fluid the momentum equation is
\begin{equation}
\rho\Big(\partial_t v + (v\!\cdot\!\nabla)v\Big) = -\nabla p + \mu\nabla^2 v,
\qquad \nabla\!\cdot v = 0 .
\end{equation}
Taking the divergence and using incompressibility eliminates the velocity time
derivative from the elliptic part and yields a Poisson equation for pressure,
\begin{equation}
\nabla^2 p = S(v) = -\rho\,\nabla\!\cdot\big((v\!\cdot\!\nabla)v\big)
+ \rho\,\nabla\!\cdot(\partial_t v),
\label{eq:ppe}
\end{equation}
whose source is built from spatial derivatives of the measured velocity. Because that
source is a product of derivatives, measurement noise (which violates
incompressibility) is amplified. A physics-informed velocity step is therefore
required first: a network fits the measured velocity while enforcing $\nabla\!\cdot v = 0$,
\begin{equation}
\mathcal{L} = \lVert v_\theta - v_\text{meas} \rVert_2^2
+ \beta\,\lVert \nabla\!\cdot v_\theta \rVert_2^2 ,
\end{equation}
producing a smooth divergence-free field whose analytic derivatives are clean. The
pressure-Poisson source and the Neumann wall flux in~\eqref{eq:ppe} are computed from
those analytic derivatives, \emph{not} by finite differences at the lumen edge, where
finite differences manufacture the worst artifacts. This analytic treatment is what
takes the recovered pressure from a non-physiological thousands of mmHg to a
physiological range~\cite{kissas2020}. The unsteady term $\partial_t v$ is
differentiated exactly by a space-time network trained over the whole cardiac cycle.

\subsection{The momentum-residual PINN failure (documented baseline)}

A single network $(x,y,z,t)\to(u,v,w,p)$ trained on the raw momentum residual
(the hidden-fluid-mechanics formulation~\cite{raissi2020}) does not recover pressure
at aortic Reynolds number: pressure is gauge-free and only weakly coupled to the
loss, so it stays near its initialization (on an analytic Poiseuille flow it recovered
under $10\%$ of the true gradient). The shipped method instead separates the
well-posed part (velocity, strongly data-constrained) from the ill-posed part
(pressure, solved by the elliptic equation). The failed approach is kept in the
repository as the documented baseline.

\subsection{Real validation with no invasive gold standard}

No non-invasive pressure gold standard exists; the validation is therefore threefold
(Table~\ref{tab:flow}). \emph{An exact analytic gate}: on a converging-duct flow whose
pressure drop is known analytically, the solve recovers 4.74 vs.\ 4.73 mmHg
(correlation 1.00), and it runs in continuous integration before any real datum is
trusted. \emph{A physiological range on the real scan}: the recovered relative pressure
spans 0.79 mmHg across the segment, small and physiological for an unobstructed aorta
(the earlier three-frame finite-difference unsteady term had inflated this to
$\approx 15$ mmHg; the analytic space-time term, gated to $dv/dt$ correlation 0.995,
removes the inflation). \emph{A clinical bracket}: the simplified-Bernoulli estimate
$4\,V_\text{max}^2 = 2.51$ mmHg from the denoised peak velocity 0.791 m/s is the same
order, exactly what an unobstructed aorta should show, while the field additionally
reveals \emph{where} the pressure varies, which a single Bernoulli number cannot.

\begin{table}[!t]
\caption{Real 4D-flow pressure reconstruction (thoracic aorta, Philips, venc
120 cm/s, 16 frames).}
\label{tab:flow}
\centering
\begin{tabular}{lr}
\toprule
Quantity & Value \\
\midrule
Lumen voxels resolved & 47{,}902 \\
Peak velocity (denoised field) & 0.791 m/s \\
Raw measured speed (noise-inflated) & $\approx 1.6$ m/s \\
Divergence reduction (denoiser) & $2.3\times$ (25.4 $\to$ 11.2 s$^{-1}$) \\
PPE relative-pressure range & 0.79 mmHg \\
Clinical Bernoulli $4V_\text{max}^2$ & 2.51 mmHg \\
Analytic gate, steady (converging duct) & corr 1.00; 4.74 vs.\ 4.73 mmHg \\
Analytic gate, unsteady (Poiseuille $dv/dt$) & corr 0.995 \\
Phase-wrap aliasing corrected & 27{,}863 samples \\
Noise-perturbation sensitivity & $<0.01$ mmHg \\
\bottomrule
\end{tabular}
\end{table}

\subsection{An honest note on uncertainty}

A deep ensemble that perturbs the measured velocity with realistic phase-contrast
noise ($5\%$ of the venc) and re-runs the whole pipeline moves the recovered pressure
by under 0.01 mmHg: the divergence-free denoiser makes the pressure essentially
insensitive to velocity measurement noise. This is a genuine robustness, but it also
means an ensemble over that noise gives a near-zero, uninformative uncertainty. The
dominant uncertainty is instead the absent invasive gold standard, the lumen
segmentation, and the unsteady-term approximation, none of which such an ensemble can
quantify. We therefore report the robustness as a scalar and deliberately do not show
a per-voxel pressure uncertainty map, which would be a misleadingly uniform near-zero
field. (This is the opposite of Case~1, where the ill-posed inverse makes the ensemble
uncertainty large and informative, so it is shown per node.)

% =====================================================================
\section{Discussion}
% =====================================================================

The two cases share one methodological spine. Each separates a well-posed part that
data constrains (the body-surface potentials; the velocity field) from an ill-posed
part that the physics forces out (the heart-surface potentials via a regularized
inverse; the pressure via an elliptic equation). Each validates against something
real: a measured gold standard where one exists (the ECGi cage), and an exact
analytic gate where one cannot (the converging-duct pressure drop). And each reports
what did not work: the BEM that does not beat the single-layer on real coarse
geometry, and the momentum-residual PINN that leaves pressure gauge-free. The
uncertainty treatment is case-appropriate rather than uniform: informative and shown
per node where the inverse is ill-posed (ECGi), and honestly suppressed to a scalar
where the denoiser makes it near-zero (4D-flow). Reporting the nulls and matching the
uncertainty to the problem is what separates a validated methodological result from a
demonstration.

% =====================================================================
\section{Limitations and conclusion}
% =====================================================================

These are validated methodological results on real experimental data, not clinically
deployed tools. ECGi is shown on a torso tank and one in-situ animal, not a patient
cohort; the 4D-flow pressure has no invasive truth to check its absolute magnitude
against, so only its analytic gate, physiological range, divergence-free denoising,
and Bernoulli bracket are claimed. Raw datasets are used under data-use agreements and
are not redistributed; only the derived reconstruction artifacts are committed, and
every number here is reproducible from them. The contribution is a demonstration that,
with a real-data-only discipline and honest nulls, physics-informed reconstruction can
recover two clinically meaningful, unmeasurable cardiac fields, ECGi heart-surface
potentials and 4D-flow aortic pressure, from the non-invasive measurements that are
tied to them by a governing equation.

% =====================================================================
\section*{Data and code availability}
% =====================================================================

Code, the derived reconstruction artifacts (the ECGi catalogue and the 4D-flow
pressure trace), and the interactive viewer are at
\url{https://github.com/fsantibanezleal/CAOS_RES_CardioPINN} under the MIT license;
a static viewer is live at \url{https://cardiopinn.fasl-work.com}. ECGi data are from
the EDGAR repository (Consortium for ECG Imaging); the 4D-flow scan is used under its
data-use agreement. Raw datasets are not redistributed.

% =====================================================================
\begin{thebibliography}{99}
\itemsep 0pt

\bibitem{raissi2019} M.~Raissi, P.~Perdikaris, and G.~E. Karniadakis,
``Physics-informed neural networks: A deep learning framework for solving forward and
inverse problems involving nonlinear partial differential equations,'' \emph{J.
Comput. Phys.}, vol.~378, pp.~686--707, 2019, doi:10.1016/j.jcp.2018.10.045.

\bibitem{raissi2020} M.~Raissi, A.~Yazdani, and G.~E. Karniadakis, ``Hidden fluid
mechanics: Learning velocity and pressure fields from flow visualizations,''
\emph{Science}, vol.~367, no.~6481, pp.~1026--1030, 2020, doi:10.1126/science.aaw4741.

\bibitem{kissas2020} G.~Kissas, Y.~Yang, E.~Hwuang, W.~R. Witschey, J.~A. Detre, and
P.~Perdikaris, ``Machine learning in cardiovascular flows modeling: Predicting arterial
blood pressure from non-invasive 4D flow MRI data using physics-informed neural
networks,'' \emph{Comput. Methods Appl. Mech. Eng.}, vol.~358, 112623, 2020,
doi:10.1016/j.cma.2019.112623.

\bibitem{krittian2012} S.~B.~S. Krittian, P.~Lamata, C.~Michler, \emph{et al.}, ``A
finite-element approach to the direct computation of relative cardiovascular pressure
from time-resolved MR velocity data,'' \emph{Med. Image Anal.}, vol.~16, no.~5,
pp.~1029--1037, 2012, doi:10.1016/j.media.2012.04.003.

\bibitem{ramanathan2004} C.~Ramanathan, R.~N. Ghanem, P.~Jia, K.~Ryu, and Y.~Rudy,
``Noninvasive electrocardiographic imaging for cardiac electrophysiology and
arrhythmia,'' \emph{Nat. Med.}, vol.~10, no.~4, pp.~422--428, 2004,
doi:10.1038/nm1011.

\bibitem{aras2015} K.~Aras, W.~Good, J.~Tate, \emph{et al.}, ``Experimental data and
geometric analysis repository (EDGAR),'' \emph{J. Electrocardiol.}, vol.~48, no.~6,
pp.~975--981, 2015, doi:10.1016/j.jelectrocard.2015.08.008.

\bibitem{cluitmans2018} M.~Cluitmans \emph{et al.}, ``Validation and opportunities of
electrocardiographic imaging: From technical achievements to clinical applications,''
\emph{Front. Physiol.}, vol.~9, 1305, 2018, doi:10.3389/fphys.2018.01305.

\bibitem{vanoosterom1983} A.~van Oosterom and J.~Strackee, ``The solid angle of a
plane triangle,'' \emph{IEEE Trans. Biomed. Eng.}, vol.~BME-30, no.~2, pp.~125--126,
1983, doi:10.1109/TBME.1983.325207.

\bibitem{hansen1992} P.~C. Hansen, ``Analysis of discrete ill-posed problems by means
of the L-curve,'' \emph{SIAM Rev.}, vol.~34, no.~4, pp.~561--580, 1992,
doi:10.1137/1034115.

\bibitem{lakshminarayanan2017} B.~Lakshminarayanan, A.~Pritzel, and C.~Blundell,
``Simple and scalable predictive uncertainty estimation using deep ensembles,'' in
\emph{Adv. Neural Inf. Process. Syst. (NeurIPS)}, 2017, arXiv:1612.01474.

\end{thebibliography}

\end{document}
